%\setcounter{MaxMatrixCols}{30}

\documentclass{article}%
\usepackage{amsmath}
\usepackage{amsfonts}
\usepackage{amssymb}
\usepackage[margin=1in]{geometry}
\usepackage{graphicx}
\usepackage[numbers]{natbib}
\usepackage{cleveref}%
\setcounter{MaxMatrixCols}{30}
%TCIDATA{OutputFilter=latex2.dll}
%TCIDATA{Version=5.50.0.2960}
%TCIDATA{CSTFile=40 LaTeX article.cst}
%TCIDATA{Created=Thursday, January 09, 2025 15:40:48}
%TCIDATA{LastRevised=Friday, March 28, 2025 09:18:44}
%TCIDATA{<META NAME="GraphicsSave" CONTENT="32">}
%TCIDATA{<META NAME="SaveForMode" CONTENT="1">}
%TCIDATA{BibliographyScheme=Manual}
%TCIDATA{<META NAME="DocumentShell" CONTENT="Standard LaTeX\Blank - Standard LaTeX Article">}
%BeginMSIPreambleData
\providecommand{\U}[1]{\protect\rule{.1in}{.1in}}
%EndMSIPreambleData
\newtheorem{theorem}{Theorem}
\newtheorem{acknowledgement}[theorem]{Acknowledgement}
\newtheorem{algorithm}[theorem]{Algorithm}
\newtheorem{axiom}[theorem]{Axiom}
\newtheorem{case}[theorem]{Case}
\newtheorem{claim}[theorem]{Claim}
\newtheorem{conclusion}[theorem]{Conclusion}
\newtheorem{condition}[theorem]{Condition}
\newtheorem{conjecture}[theorem]{Conjecture}
\newtheorem{corollary}[theorem]{Corollary}
\newtheorem{criterion}[theorem]{Criterion}
\newtheorem{definition}[theorem]{Definition}
\newtheorem{example}[theorem]{Example}
\newtheorem{exercise}[theorem]{Exercise}
\newtheorem{lemma}[theorem]{Lemma}
\newtheorem{notation}[theorem]{Notation}
\newtheorem{problem}[theorem]{Problem}
\newtheorem{proposition}[theorem]{Proposition}
\newtheorem{remark}[theorem]{Remark}
\newtheorem{solution}[theorem]{Solution}
\newtheorem{summary}[theorem]{Summary}
\newenvironment{proof}[1][Proof]{\noindent\textbf{#1.} }{\ \rule{0.5em}{0.5em}}



\DeclareMathOperator{\supp}{supp}
\newcommand{\R}{\mathbb{R}}
\newcommand{\E}{\mathbb{E}}
\newcommand{\Prob}{\mathbb{P}}
\newcommand{\dd}{\,\mathrm{d}}
\newcommand{\thetahat}{\hat{\theta}_n}
\newcommand{\thetastar}{\theta^*}
\newcommand{\one}{\boldsymbol{1}}



\begin{document}

\title{Bayes-assisted anytime-valid t-tests and confidence sequences for Gaussian means with unknown variance}
\author{Fran\c{c}ois}
\maketitle

\section{Calculations}



% compute_An
% numerical integration via integrate.quad
% compute_mixture_martingale
% numerical integration via integrate.quad
% compute_g_inverse_lebesgue
% numerical root-finding via optimize.brentq
% compute_calibration_function
% Monte Carlo simulation
% compute_g_inverse
% numerical root-finding, repeatedly calling Monte Carlo calibration
% compute_ville_pvalue
% depends on numerical compute_mixture_martingale
% compute_extended_pvalue_lebesgue
% closed-form once M_n^{Leb}
% is available, so not really numerical in itself
% compute_extended_pvalue
% numerical root-finding plus numerical integration plus Monte Carlo
% experiment2_pvalue_trajectories
% simulation experiment with repeated numerical evaluations



\[
H_0:\theta^* = 0
\qquad\text{vs}\qquad
H_1:\theta^*\neq 0,
\]

where \(Y_i \sim N(\mu,\sigma^2)\), both \(\mu,\sigma\) unknown, and

\[
\theta^*=\frac{\mu}{\sigma}.
\]

The scale-invariant statistic is defined as

\[
u_n
=
\frac{\sum_{i=1}^n Y_i}{\sqrt{\sum_{i=1}^n Y_i^2}}
=
\frac{n\hat\theta_n}{\sqrt{(n-1)+n\hat\theta_n^2}},
\]

and the building block

\[
h_{\theta,n}(u)
=
e^{-n\theta^2/2}\,A_n(\theta u),
\]

with

\[
A_n(z)
=
\frac{1}{\Gamma(n/2)}
\int_0^\infty y^{n/2-1}e^{-y+z\sqrt{2y}}\,dy
=
\mathbb E[e^{zR_n}],
\qquad R_n\sim \chi_n.
\]

Under \(H_0:\theta^*=0\), one has \(h_{0,n}(u)=1\), so the likelihood-ratio martingale reduces to \(L_n(y_{1:n};0,\theta)=h_{\theta,n}(u_n)\).

\paragraph{General mixture martingale:}

Given a prior or \(\sigma\)-finite mixing measure \(\Pi\) on \(\mathbb R\), the mixture martingale is

\[
M_n^\Pi(u)
=
\int_{\mathbb R} h_{\theta,n}(u)\,\Pi(d\theta)
=
\int_{\mathbb R} e^{-n\theta^2/2}A_n(\theta u)\,\Pi(d\theta).
\]

If \(\Pi\) has density \(\pi(\theta)\), this becomes

\[
M_n^\Pi(u)
=
\int_{-\infty}^\infty e^{-n\theta^2/2}A_n(\theta u)\,\pi(\theta)\,d\theta.
\]


Starting from the mixture construction in the document:

\[
M_n(y_{1:n};\theta_0)
=
\int_{\mathbb R}L_n(y_{1:n};\theta_0,\theta)\,\Pi(d\theta)
=
\frac{\int h_{\theta,n}(y_{1:n})\,\Pi(d\theta)}{h_{\theta_0,n}(y_{1:n})}.
\]

For \(\theta_0=0\), \(h_{0,n}=1\), hence

\[
M_n^\Pi(y_{1:n})
=
\int h_{\theta,n}(y_{1:n})\,\Pi(d\theta)
=
\int e^{-n\theta^2/2}A_n(\theta u_n)\,\Pi(d\theta).
\]


\paragraph{Calibration function \(g_\Pi\):}

For a proper prior, the Ville p-process is

\[
P_{n}^{V,\Pi}
=
\frac{1}{M_n^\Pi}.
\]

For a general \(\sigma\)-finite \(\Pi\), the document uses the extended calibration

\[
g_\Pi(c)
=
\mathbb E_0\!\left[\min\!\left(\frac{M_2^\Pi(U_2)}{c},1\right)\right],
\]

where the expectation is under \(H_0:\theta^*=0\), and \(U_2\) is the 2-sample version of \(u_n\). The extended p-variable is then

\[
P_n^{E,\Pi}
=
g_\Pi(M_n^\Pi).
\]

%This is the mathematical definition behind `compute_calibration_function` and `compute_extended_pvalue`.


We take \(n=2\) in \(g_\Pi\) because the extended Ville calibration is defined from \(M_2\), and for the improper Lebesgue prior the process naturally starts at \(n=2\).



\paragraph{Prior-by-prior formulas:}


\begin{itemize}
    \item[A.] \textbf{Gaussian prior:}
    
    Suppose

\[
\Pi(d\theta)=\pi_G(\theta)\,d\theta,
\qquad
\pi_G(\theta)
=
\frac{1}{\sqrt{2\pi}s}
\exp\!\left(-\frac{(\theta-m)^2}{2s^2}\right).
\]

Then

\[
M_n^{G}(u)
=
\int_{-\infty}^\infty
e^{-n\theta^2/2}A_n(\theta u)
\frac{1}{\sqrt{2\pi}s}
\exp\!\left(-\frac{(\theta-m)^2}{2s^2}\right)
\,d\theta.
\]

The calibration is

\[
g_G(c)
=
\mathbb E_0\!\left[
\min\!\left(\frac{M_2^G(U_2)}{c},1\right)
\right].
\]

Substituting the Gaussian density into the general mixture formula:

\[
M_n^\Pi(u)=\int h_{\theta,n}(u)\pi(\theta)\,d\theta.
\]

There is no closed form given for \(M_n^G\) or \(g_G\), so the code uses numerical integration for \(M_n^G\) and Monte Carlo for \(g_G\).



\item[B.] \textbf{Laplace prior:}

Let

\[
\pi_L(\theta)
=
\frac{1}{2b}\exp\!\left(-\frac{|\theta-m|}{b}\right).
\]

Then

\[
M_n^{L}(u)
=
\int_{-\infty}^\infty
e^{-n\theta^2/2}A_n(\theta u)
\frac{1}{2b}\exp\!\left(-\frac{|\theta-m|}{b}\right)\,d\theta,
\]

and

\[
g_L(c)
=
\mathbb E_0\!\left[
\min\!\left(\frac{M_2^L(U_2)}{c},1\right)
\right].
\]

Again, this is just substitution into the general formula.


\item[C.] \textbf{Student-\(t\) prior:}

For \(\nu>0\), center \(m\), scale \(s\),

\[
\pi_T(\theta)
=
\frac{1}{s}
\frac{\Gamma\!\left(\frac{\nu+1}{2}\right)}
{\sqrt{\nu\pi}\,\Gamma\!\left(\frac{\nu}{2}\right)}
\left(
1+\frac{1}{\nu}\left(\frac{\theta-m}{s}\right)^2
\right)^{-(\nu+1)/2}.
\]

Then

\[
M_n^T(u)
=
\int_{-\infty}^\infty
e^{-n\theta^2/2}A_n(\theta u)\,\pi_T(\theta)\,d\theta,
\]

and

\[
g_T(c)
=
\mathbb E_0\!\left[
\min\!\left(\frac{M_2^T(U_2)}{c},1\right)
\right].
\]

The document studies shifted polynomial-tail priors \(\pi_m(\theta)=\pi(\theta-m)\), and Student-\(t\) is one of the main examples with tail index

\[
\pi(x)\sim C|x|^{-(\nu+1)}.
\]

This is why the asymptotic recovery of the Lebesgue prior applies here.


\item[D.] \textbf{Cauchy prior:}

This is Student-\(t\) with \(\nu=1\):

\[
\pi_C(\theta)
=
\frac{1}{\pi s}\frac{1}{1+\left(\frac{\theta-m}{s}\right)^2}.
\]

So

\[
M_n^C(u)
=
\int_{-\infty}^\infty
e^{-n\theta^2/2}A_n(\theta u)\,
\frac{1}{\pi s}\frac{1}{1+\left(\frac{\theta-m}{s}\right)^2}
\,d\theta,
\]

and

\[
g_C(c)
=
\mathbb E_0\!\left[
\min\!\left(\frac{M_2^C(U_2)}{c},1\right)
\right].
\]

The document explicitly mentions the Cauchy prior as a main example.


\item[E.] \textbf{Horseshoe prior:}

We use an approximation

\[
\pi_{HS}(\theta)
\approx
\frac{1}{\pi s}
\frac{\log\!\left(1+\frac{1}{z^2}\right)}{1+z^2},
\qquad
z=\frac{\theta-m}{s},
\]

with a small-\(z\) regularization near the origin ($\approx \frac{1}{\pi s}$).

So the corresponding mixture is

\[
M_n^{HS}(u)
=
\int_{-\infty}^\infty
e^{-n\theta^2/2}A_n(\theta u)\,\pi_{HS}(\theta)\,d\theta,
\]

and

\[
g_{HS}(c)
=
\mathbb E_0\!\left[
\min\!\left(\frac{M_2^{HS}(U_2)}{c},1\right)
\right].
\]

The horseshoe location family has polynomial tail index \(\beta=2\), so it also falls under the same “recovery of Lebesgue under large shift” theorem.

\end{itemize}

\paragraph{Improper Lebesgue prior:}

Take

\[
\Pi_{\mathrm{Leb}}(d\theta)=d\theta.
\]

Then

\[
M_n^{\mathrm{Leb}}(u)
=
\int_{-\infty}^\infty e^{-n\theta^2/2}A_n(\theta u)\,d\theta.
\]

We have the closed form

\[
M_n^{\mathrm{Leb}}(u)
=
\sqrt{\frac{2\pi}{n}}
\left(1-\frac{u^2}{n}\right)^{-n/2},
\qquad |u|<\sqrt n.
\]

%That is exactly your `compute_mixture_lebesgue`.

% ### Derivation
% Use

% \[
% A_n(\theta u)=\mathbb E[e^{\theta u R_n}],
% \qquad R_n\sim \chi_n.
% \]

% Then by Tonelli,

% \[
% M_n^{\mathrm{Leb}}(u)
% =
% \int_{\mathbb R} e^{-n\theta^2/2}\mathbb E[e^{\theta u R_n}]\,d\theta
% =
% \mathbb E\!\left[
% \int_{\mathbb R}\exp\!\left(-\frac n2\theta^2+\theta uR_n\right)\,d\theta
% \right].
% \]

% Complete the square:

% \[
% -\frac n2\theta^2+\theta uR_n
% =
% -\frac n2\left(\theta-\frac{uR_n}{n}\right)^2
% +\frac{u^2R_n^2}{2n}.
% \]

% Hence

% \[
% \int_{\mathbb R}\exp\!\left(-\frac n2\theta^2+\theta uR_n\right)\,d\theta
% =
% \sqrt{\frac{2\pi}{n}}\exp\!\left(\frac{u^2R_n^2}{2n}\right).
% \]

% Therefore

% \[
% M_n^{\mathrm{Leb}}(u)
% =
% \sqrt{\frac{2\pi}{n}}
% \mathbb E\!\left[\exp\!\left(\frac{u^2R_n^2}{2n}\right)\right].
% \]

% Since \(R_n^2\sim \chi_n^2\), its mgf is

% \[
% \mathbb E[e^{tR_n^2}] = (1-2t)^{-n/2},
% \qquad t<1/2.
% \]

% Take \(t=u^2/(2n)\). Then

% \[
% M_n^{\mathrm{Leb}}(u)
% =
% \sqrt{\frac{2\pi}{n}}
% \left(1-\frac{u^2}{n}\right)^{-n/2}.
% \]

% This is Proposition 2.



\textbf{Lebesgue calibration \(g_{\mathrm{Leb}}\)}

We also have the closed form for

\[
g_{\mathrm{Leb}}(c)
=
\mathbb E_0\!\left[
\min\!\left(\frac{M_2^{\mathrm{Leb}}(U_2)}{c},1\right)
\right].
\]

It is

\[
g_{\mathrm{Leb}}(c)
=
\begin{cases}
1, & c\le \sqrt{\pi},\\[1ex]
\displaystyle
\frac{2}{\pi}
\left[
\arcsin\!\left(\sqrt{\frac{\sqrt\pi}{c}}\right)
+
\sqrt{\frac{\sqrt\pi}{c}\left(1-\frac{\sqrt\pi}{c}\right)}
\right],
& c>\sqrt{\pi}.
\end{cases}
\]

%That is exactly your `compute_g_lebesgue`.

\textbf{Derivation:}


For \(n=2\),

\[
M_2^{\mathrm{Leb}}(u)=\sqrt{\pi}\left(1-\frac{u^2}{2}\right)^{-1}.
\]

Under \(H_0\), define

\[
B:=1-\frac{U_2^2}{2}.
\]

The document shows

\[
B\sim \mathrm{Beta}\!\left(\frac12,\frac12\right),
\qquad
f_B(b)=\frac{1}{\pi\sqrt{b(1-b)}},\quad b\in(0,1),
\]

and

\[
M_2^{\mathrm{Leb}}(U_2)=\frac{\sqrt{\pi}}{B}.
\]

Hence

\[
g_{\mathrm{Leb}}(c)
=
\int_0^1
\min\!\left(\frac{\sqrt{\pi}}{cb},1\right)
\frac{db}{\pi\sqrt{b(1-b)}}.
\]

If \(c\le \sqrt\pi\), the fraction is always at least \(1\), so \(g_{\mathrm{Leb}}(c)=1\).

If \(c>\sqrt\pi\), let

\[
a=\frac{\sqrt\pi}{c}\in(0,1).
\]

Then split the integral at \(b=a\):

\[
g_{\mathrm{Leb}}(c)
=
\int_0^a \frac{db}{\pi\sqrt{b(1-b)}}
+
a\int_a^1 \frac{db}{\pi b\sqrt{b(1-b)}}.
\]

Evaluating these gives the closed form above. That is Proposition 3.


\paragraph{Shifted polynomial-tail priors:}

We have a shifted prior

\[
\pi_m(\theta)=\pi(\theta-m),
\]

and then the corresponding mixture

\[
M_{n,m}(u)
=
\int_{\mathbb R} e^{-n\theta^2/2}A_n(\theta u)\pi(\theta-m)\,d\theta.
\]

It also defines

\[
a_m := \pi(-m),
\]

and proves that for polynomial-tail priors,

\[
\frac{M_{n,m}(u)}{a_m}\to M_n^{\mathrm{Leb}}(u)
\qquad (m\to\infty),
\]

and similarly

\[
g_m(a_m c)\to g_{\mathrm{Leb}}(c).
\]

This is the reason the shifted Student-\(t\), Cauchy, and horseshoe priors recover the improper Lebesgue benchmark under extreme prior-data conflict.


% \paragraph{Improper exponential prior \(e^{-\kappa\theta}\):}

% The code discussed the improper prior

% \[
% \Pi_\kappa(d\theta)=e^{-\kappa|\theta|}\,d\theta,
% \qquad \kappa\ge 0.
% \]

% For \(\kappa=0\), this is the Lebesgue prior.

% The corresponding mixture is

% \[
% M_n^{(\kappa)}(u)
% =
% \int_{\mathbb R}
% e^{-n\theta^2/2}A_n(\theta u)e^{-\kappa|\theta|}
% \,d\theta.
% \]

% The document gives the closed-form representation

% \[
% M_n^{(\kappa)}(u)
% =
% \sqrt{\frac{2\pi}{n}}\,
% \mathbb E\!\left[
% \exp\!\left(\frac{(uR_n-\kappa)^2}{2n}\right)
% \right].
% \]

% For \(\kappa=0\), this reduces to the Lebesgue formula.

% \textbf{Derivation:}

% Start from

% \[
% M_n^{(\kappa)}(u)
% =
% \int_{\mathbb R} e^{-n\theta^2/2}A_n(\theta u)e^{-\kappa|\theta|}\,d\theta.
% \]

% The document’s derivation effectively rewrites the exponent in linear form inside the expectation and then completes the square, obtaining

% \[
% M_n^{(\kappa)}(u)
% =
% \mathbb E\!\left[
% \int_{\mathbb R}
% \exp\!\left(-\frac{n\theta^2}{2}+\theta(uR_n-\kappa)\right)d\theta
% \right].
% \]

% Complete the square:

% \[
% -\frac{n\theta^2}{2}+\theta(uR_n-\kappa)
% =
% -\frac n2\left(\theta-\frac{uR_n-\kappa}{n}\right)^2
% +\frac{(uR_n-\kappa)^2}{2n}.
% \]

% Then the Gaussian integral contributes \(\sqrt{2\pi/n}\), giving

% \[
% M_n^{(\kappa)}(u)
% =
% \sqrt{\frac{2\pi}{n}}\,
% \mathbb E\!\left[
% \exp\!\left(\frac{(uR_n-\kappa)^2}{2n}\right)
% \right].
% \]



% \paragraph{Calibration for the exponential improper prior:}


% Its calibration is defined exactly the same way:

% \[
% \tilde g_\kappa(c)
% =
% \mathbb E_0\!\left[
% \min\!\left(\frac{M_2^{(\kappa)}(U_2)}{c},1\right)
% \right].
% \]

% For \(n=2\),

% \[
% M_2^{(\kappa)}(u)
% =
% \sqrt{\pi}\,
% \mathbb E\!\left[
% \exp\!\left(\frac{(uR_2-\kappa)^2}{4}\right)
% \right].
% \]

% There is no closed form analogous to \(g_{\mathrm{Leb}}\) known here, so one would approximate \(\tilde g_\kappa\) numerically or by Monte Carlo.



% \section{Summary table}

% For all priors, the common structure is

% \[
% M_n^\Pi(u)=\int e^{-n\theta^2/2}A_n(\theta u)\,\Pi(d\theta),
% \qquad
% g_\Pi(c)=\mathbb E_0\!\left[\min\!\left(\frac{M_2^\Pi(U_2)}{c},1\right)\right].
% \]

% Then:

% Gaussian
% \[
% \Pi(d\theta)=\pi_G(\theta)\,d\theta
% \]
% \[
% M_n^G(u)=\int e^{-n\theta^2/2}A_n(\theta u)\pi_G(\theta)\,d\theta
% \]
% \[
% g_G(c)=\mathbb E_0[\min(M_2^G(U_2)/c,1)]
% \]

% - **Laplace**
% \[
% \Pi(d\theta)=\pi_L(\theta)\,d\theta
% \]
% \[
% M_n^L(u)=\int e^{-n\theta^2/2}A_n(\theta u)\pi_L(\theta)\,d\theta
% \]

% - **Student-\(t\)**
% \[
% \Pi(d\theta)=\pi_T(\theta)\,d\theta
% \]
% \[
% M_n^T(u)=\int e^{-n\theta^2/2}A_n(\theta u)\pi_T(\theta)\,d\theta
% \]

% - **Cauchy**
% \[
% \Pi(d\theta)=\pi_C(\theta)\,d\theta
% \]
% \[
% M_n^C(u)=\int e^{-n\theta^2/2}A_n(\theta u)\pi_C(\theta)\,d\theta
% \]

% - **Horseshoe approximation**
% \[
% \Pi(d\theta)=\pi_{HS}(\theta)\,d\theta
% \]
% \[
% M_n^{HS}(u)=\int e^{-n\theta^2/2}A_n(\theta u)\pi_{HS}(\theta)\,d\theta
% \]

% - **Lebesgue**
% \[
% \Pi(d\theta)=d\theta
% \]
% \[
% M_n^{\mathrm{Leb}}(u)=\sqrt{\frac{2\pi}{n}}\left(1-\frac{u^2}{n}\right)^{-n/2}
% \]
% \[
% g_{\mathrm{Leb}}(c)=
% \begin{cases}
% 1, & c\le \sqrt\pi,\\[1ex]
% \frac{2}{\pi}\left[\arcsin\!\left(\sqrt{\frac{\sqrt\pi}{c}}\right)+
% \sqrt{\frac{\sqrt\pi}{c}\left(1-\frac{\sqrt\pi}{c}\right)}\right], & c>\sqrt\pi
% \end{cases}
% \]

% - **Improper exponential**
% \[
% \Pi_\kappa(d\theta)=e^{-\kappa|\theta|}d\theta
% \]
% \[
% M_n^{(\kappa)}(u)
% =
% \sqrt{\frac{2\pi}{n}}\,
% \mathbb E\!\left[\exp\!\left(\frac{(uR_n-\kappa)^2}{2n}\right)\right]
% \]
% \[
% \tilde g_\kappa(c)
% =
% \mathbb E_0\!\left[\min(M_2^{(\kappa)}(U_2)/c,1)\right]
% \]

% ---

% # 8. Relation to your code

% Your code is implementing these formulas as follows:

% - `compute_mixture_martingale`:
% \[
% M_n^\Pi(u)=\int h_{\theta,n}(u)\pi(\theta)\,d\theta
% \]

% - `compute_mixture_lebesgue`:
% closed form for \(M_n^{\mathrm{Leb}}(u)\)

% - `compute_g_lebesgue`:
% closed form for \(g_{\mathrm{Leb}}(c)\)

% - `compute_calibration_function`:
% Monte Carlo approximation of
% \[
% g_\Pi(c)=\mathbb E_0[\min(M_2^\Pi(U_2)/c,1)]
% \]

% - `compute_ville_pvalue`:
% \[
% P_n^{V,\Pi}=\frac{1}{M_n^\Pi}
% \]

% - `compute_extended_pvalue_lebesgue`:
% \[
% P_n^{E,\mathrm{Leb}}=g_{\mathrm{Leb}}(M_n^{\mathrm{Leb}})
% \]

% - `compute_extended_pvalue`:
% numerical inversion-based approximation to
% \[
% P_n^{E,\Pi}=g_\Pi(M_n^\Pi).
% \]

% If you want, I can next turn this into a one-page “notation sheet” with every formula written in the exact symbols you want for your report.






% \bibliography{bayesanytimettest}
% \bibliographystyle{alpha}

\end{document} 